\section{Curious hard Lefschetz on the full character variety}
\label{sec:betti}

We now establish the Betti Lefschetz statement needed for the comparison
criterion.  The point requiring care is the full cohomology of the
fixed-determinant space.  We obtain it from a finite cover of a
$\mathrm{GL}_n$ family, retaining every component of every lifted stratum.

For a finite-dimensional vector space graded by half weights, a class of
cohomological degree two and weight four acts with grading degree two.
We say it satisfies \emph{Lefschetz with center $c$} if its $i$th power
is an isomorphism from grade $c-i$ to grade $c+i$, in the corresponding
cohomological degrees, for every $i\geq0$.

A mixed Hodge structure is \emph{Hodge--Tate} if its odd weight-graded
pieces vanish and its weight-$2r$ piece has only Hodge type $(r,r)$.

\begin{theorem}\label{betti:lefschetz}
The mixed Hodge structure on $H^*(X^{\mathrm B},\mathbb Q)$ is Hodge--Tate.
Under the nonabelian-Hodge identification, the class $e$ has nonzero image
in $\operatorname{Gr}^W_4H^2(X^{\mathrm B},\mathbb Q)$, and
\[
 e^i:\operatorname{Gr}^W_{2R-2i}H^m(X^{\mathrm B},\mathbb Q)
 \xrightarrow{\ \sim\ }
 \operatorname{Gr}^W_{2R+2i}H^{m+2i}(X^{\mathrm B},\mathbb Q)
 \qquad(i\geq0).
\]
In particular, $W_{2k}=W_{2k+1}$ on this cohomology.  These assertions
concern the entire cohomology, including all nontrivial
$\Gamma$-character summands after complexification.
\end{theorem}

\subsection{The ordered family and its scalar fiber}

Let
\[
 V=\left\{(t_1,\ldots,t_n)\in(\mathbb C^*)^n:
 \prod_i t_i=1,\quad
 \prod_{i\in I}t_i\ne1\text{ for }\varnothing\ne I\subsetneq\{1,\ldots,n\}
 \right\},
 \qquad q=(\zeta,\ldots,\zeta).
\]
The space $V$ is a connected smooth open subset of a torus, and $q\in V$
since $\zeta$ has order $n$.  Write $\nu=\prod_{j=1}^g[A_j,B_j]$.
The fiber $Y^{\mathrm B}_t$ parametrizes tuples
$(A_1,B_1,\ldots,A_g,B_g)\in\mathrm{SL}_n^{2g}$ together with a full flag
preserved by $\nu$, on whose ordered graded lines $\nu$ acts by
$t_1,\ldots,t_n$, modulo simultaneous conjugation.  Repeated eigenvalues
are allowed.

Every such tuple is irreducible.  Indeed, on a common invariant proper
subspace the determinant of $\nu$ is one, while its eigenvalues form a
proper nonempty submultiset of the $t_i$.  This contradicts the definition
of $V$.  Thus conjugation gives geometric quotients by a free
$\mathrm{PGL}_n$ action.  They may be constructed by the usual
\'etale-local principal-bundle charts on irreducible tuples; with a full
flag, these charts can also be obtained by completing its first line to
a basis using words in the matrices.

\begin{lemma}\label{betti:geometry}
The morphism $Y^{\mathrm B}\to V$ is smooth of relative dimension
$2R+2b$.  There is a correspondence
\[
 X^{\mathrm B}\xleftarrow{\ \pi\ }Z^{\mathrm B}
 \xrightarrow{\ \iota\ }Y^{\mathrm B}_q,
\]
where $Z^{\mathrm B}$ imposes $\nu=\zeta I$ and $\pi$ is the full flag
bundle.  The map $\iota$ is a regular embedding of codimension $b$ with
normal bundle $T_\pi^*$.  The cohomological operators
\[
 \mathsf A=\pi_*\iota^*,\qquad
 \mathsf B=\iota_*\pi^*
\]
have degrees $-2b$ and $2b$, respectively, and satisfy
\begin{equation}\label{betti:split}
 \mathsf A\mathsf B=(-1)^b n!\,\mathrm{id}.
\end{equation}
Here and below a covariant star denotes the Gysin map.
\end{lemma}

\begin{proof}
The commutator-product map is a submersion at every irreducible tuple.
We give the differential check.  For a pair $a,b'$ set
$c=ab'a^{-1}b'^{-1}$.  Its right-translated differential has terms
\[
 (1-\operatorname{Ad}_{ab'a^{-1}})u,\qquad
 (\operatorname{Ad}_a-\operatorname{Ad}_c)v.
\]
Using the trace pairing, an annihilating vector $z$ is fixed by
$ab'a^{-1}$ and satisfies
$\operatorname{Ad}_{a^{-1}}z=\operatorname{Ad}_{b'}z$.
It is therefore fixed by $ab'$, hence by $a$ and $b'$.
For a product of commutators one applies this argument to the testing
vector conjugated by each prefix.  Since the vector is fixed by the
current commutator, it stays unchanged from one pair to the next.
It consequently centralizes every matrix.  Irreducibility forces a
scalar, and the trace-free condition forces zero.

The space of a matrix in a flag Borel, with its flag, is smooth over
its ordered diagonal torus, with fibers of dimension $2b$.
Pulling this space back along the submersion and then taking the free
quotient gives smoothness.  The relative dimension is
\[
 2g(n^2-1)+2b-2(n^2-1)=2R+2b.
\]
At $q$, the section $\zeta^{-1}\nu-I$ takes values in the nilradical
bundle over the flag space.  Its zero locus is $Z^{\mathrm B}$, and
submersivity proves that this section is regular.  The trace pairing
identifies the nilradical with the dual of the flag tangent space, so
its associated bundle is $T_\pi^*$.  Self-intersection and the projection
formula now give
\[
 \mathsf A\mathsf B(\alpha)
 =\pi_*\bigl(c_b(T_\pi^*)\pi^*\alpha\bigr)
 =(-1)^b\chi(\mathrm{Fl}_n)\alpha
 =(-1)^b n!\alpha.
\]
All constructions descend through $\mathrm{PGL}_n$ and apply to full
cohomology.
\end{proof}

\subsection{Lifting the toric stratifications}

Set $D_1=(\mathbb C^*)^{2g}$ and
$\widehat Y^{\mathrm B}=Y^{\mathrm B}\times D_1$.
Multiplying each of the $2g$ matrices by its corresponding scalar gives
\begin{equation}\label{betti:cover}
 \widehat Y^{\mathrm B}\longrightarrow Y^{\mathrm B}_{\mathrm{GL}},
\end{equation}
where the target is the same ordered family with matrices in
$\mathrm{GL}_n$.  This is the pullback of the $n$th-power map on the
$2g$ determinant coordinates, hence is finite \'etale of degree
$n^{2g}$.  Its source is the indicated product: dividing by the chosen
determinant roots gives matrices in $\mathrm{SL}_n$, without changing
commutators or flags.

We use Mellit's toric stratifications of the $\mathrm{GL}_n$ family
\cite{Mellit2025}.  To identify the family exactly, take two punctures
with first monodromy $\zeta^{-1}I$ and last monodromy
$\zeta\nu^{-1}$.  The latter has ordered eigenvalues
$(\zeta t_i^{-1})_i$.  The product relation is
$\nu\zeta^{-1}I\cdot\zeta\nu^{-1}=I$, and the genericity condition of
\cite[Definition~4.9]{Mellit2025} is precisely the subproduct condition
in the definition of $V$.  Its identity last monodromy corresponds to
$q$.  The ordering condition in that reference is automatic for the
scalar first monodromy and for the distinct-eigenvalue last monodromy.

We recall the two features of the stratifications which will be used.
First, at a parameter with distinct $t_i$, an auxiliary rank-$b$ vector
bundle over $Y^{\mathrm B}_{\mathrm{GL},t}$ has a finite stratification
refined into vector bundles over products of tori and affine spaces.
The restriction of
Mellit's degree-two class gives compact-support curious Lefschetz on
every final stratum, with the same middle weight.  If
$\delta=\dim Y^{\mathrm B}_{\mathrm{GL},t}$, that middle \emph{actual}
weight is $\delta+2b$; removal of the vector-bundle shift gives
$\delta$ \cite[Theorem~6.14, Section~7.4, and Theorem~7.5]{Mellit2025}.
Second, over all of $V$, an auxiliary affine bundle of rank $b$ has a
finite stratification whose refinements are vector bundles over
\[
 \mathbb A^a\times(\mathbb C^*)^s\times V
\]
with structural map the projection to $V$
\cite[proof of Proposition~8.7]{Mellit2025}.
These assertions are about the final torus--affine refinements, not
about arbitrary braid strata occurring earlier in their construction.

Here is the additional finite-cover argument needed for
\eqref{betti:cover}.

\begin{lemma}\label{betti:finite-cover}
Pulling the preceding auxiliary bundles and stratifications back along
\eqref{betti:cover} has the following consequences.
\begin{enumerate}
\item Every fiber of $\widehat Y^{\mathrm B}\to V$ and of
$Y^{\mathrm B}\to V$ has Hodge--Tate cohomology.  For
$\widehat Y^{\mathrm B}$ at a parameter with distinct eigenvalues, the pulled-back
class satisfies curious hard Lefschetz on ordinary cohomology with
half-weight center $\delta/2=R+b+g$.
\item The cohomology sheaves of both ordinary and compact-support
integral direct images over $V$ are locally constant and finitely
generated.  Ordinary direct image commutes with restriction to every
fiber.  The same assertions hold for $Y^{\mathrm B}\to V$.
\end{enumerate}
\end{lemma}

\begin{proof}
Consider first a final stratum $F=(\mathbb C^*)^s\times\mathbb A^a$,
or a vector bundle over such a stratum.  Its homotopy type is that of
a torus.  Each connected component $\widetilde F$ of its finite \'etale
preimage corresponds to a finite-index subgroup of $\mathbb Z^s$.
The pullback
\[
 H^*(F,\mathbb Q)\longrightarrow H^*(\widetilde F,\mathbb Q)
\]
is therefore an isomorphism: both rings are the exterior algebras of
the rational duals of these lattices.  Since the map is algebraic, it
is an isomorphism of mixed Hodge structures.  Both spaces are smooth
of the same dimension, and finite-degree trace together with
Poincar\'e duality gives the same assertion for compact supports.
These isomorphisms commute with cup product by the pulled-back class.

We apply this argument to each connected component separately.
The cohomology of a disconnected preimage is their direct sum, and
all components have the same Lefschetz center.  Preimages of the closed
filtration of the original stratification still give a finite closed
filtration.  In the compact-support localization sequences, strictness
of mixed Hodge morphisms \cite[Theorem~2.3.5]{Deligne1971} permits
passage to weight-gradeds.  Comparing
these exact sequences by a fixed power of the common class and using
the five lemma passes the Lefschetz isomorphisms from the strata to
the entire auxiliary space.  The same sequences show that its
cohomology is Hodge--Tate.

The pulled-back auxiliary bundle has the same rank $b$.  Its Thom
isomorphism shifts cohomological degree and actual weight by $2b$.
Thus the compact-support middle weight drops from $\delta+2b$ to
$\delta$.  Smooth Poincar\'e duality sends weight $w$ to
$2\delta-w$, so ordinary cohomology has the same middle weight
$\delta$, or center $\delta/2$ in half weights.  The relative product
stratification proves the Hodge--Tate assertion on every fiber,
including repeated eigenvalues, by the same componentwise argument
and the affine-bundle Thom isomorphism.

For integral local constancy, use the relative stratification and
choose a small contractible analytic open set $U\subset V$.
A finite cover of $F\times U$ is pulled back from a finite cover of
$F$, since $\pi_1(F\times U)=\pi_1(F)$.  The same statement holds for
the total space of a vector bundle over $F\times U$, which is homotopy
equivalent to its base.  Topologically the bundle may be identified,
in the $U$ direction, with the pullback of its restriction over one
point.  Equivalently, one may use the integral Thom isomorphism for
its complex orientation.  It follows that the compact-support
cohomology sheaves of each covered stratum are locally constant
finitely generated abelian groups on $U$.  Localization triangles
pass this conclusion through the finite stratification: kernels,
cokernels, and extensions of locally constant sheaves are locally
constant.  The auxiliary affine-bundle trace removes its shift.

Write $f$ for the covered smooth family of relative dimension
$\delta$.  The complex $K=Rf_!\mathbb Z$ is locally perfect: its
cohomology is bounded and finitely generated, and $\mathbb Z$ has
finite global dimension.  Smooth Verdier duality gives
\[
 Rf_*\mathbb Z\simeq
 R\mathcal Hom(K,\mathbb Z_V)[-2\delta].
\]
The right side is locally constant and perfect, so taking its stalk
commutes with this duality.  Compact-support base change and the
fiberwise cup--trace pairing identify its stalk with
$R\Gamma(f^{-1}(t),\mathbb Z)$.  This identification is precisely the
natural ordinary base-change morphism: both are adjoint to the cup
product followed by the smooth orientation trace, whose formation
commutes with compact-support base change.  This proves the ordinary
assertion, without using a nonproper smooth base-change principle.
Finally, write $p:Y^{\mathrm B}\times D_1\to Y^{\mathrm B}$ for the
projection and $s$ for its section at the identity of $D_1$.  For ordinary
direct images, $s^*p^*=\mathrm{id}$ gives the required summand.  For
compact supports use the Gysin map of the closed section, followed by
integration over $D_1$: these maps have degrees $4g$ and $-4g$ and
compose to the identity.  Equivalently, the integral compact-support
K\"unneth complex of $D_1$ has its top class as a split summand.  Thus
both local-constancy assertions pass to $Y^{\mathrm B}$, and natural
ordinary base change passes through the first splitting.  The ordinary
splitting is algebraic on each fiber, so it also proves the asserted
Hodge--Tate property for $Y^{\mathrm B}$.
\end{proof}

\subsection{The scalar summand and the determinant factor}

We have now obtained Lefschetz on a regular-semisimple fiber of the
cover and cohomological local constancy over the entire ordered base.
We next pass to $q$ and then remove the flag and determinant factors.

For a smooth algebraic family whose ordinary direct images are lisse (that is, have locally constant
cohomology sheaves) and satisfy fiberwise base change, mixed Hodge module direct image
gives admissible graded-polarizable variations of mixed Hodge
structure on these local systems.  In particular, parallel transport
preserves the weight filtration; cup products are parallel as well.
The construction and its diagram version are given in
Proposition~\ref{diag:variation}.  Apply this fact to the family in
Lemma~\ref{betti:finite-cover}.  Along a path in the connected space
$V$, transport a Lefschetz class from a distinct-eigenvalue fiber to
$q$.  We obtain
\[
 \omega\in\operatorname{Gr}^W_4
 H^2(\widehat Y^{\mathrm B}_q,\mathbb C)
\]
with Lefschetz center $R+b+g$.  A choice of path suffices; no assertion
that this class is invariant under all monodromy is needed.

Take the product of the correspondence in
Lemma~\ref{betti:geometry} with $D_1$, writing
$\widehat X^{\mathrm B}=X^{\mathrm B}\times D_1$ and
$\widehat Z^{\mathrm B}=Z^{\mathrm B}\times D_1$.
The flag bundle has rational algebraic Leray--Hirsch classes.  More
explicitly, the ratios of the ordered graded lines descend for the
universal projective bundle, and their first Chern classes restrict
to generators of the flag cohomology over $\mathbb Q$.  Indeed,
differences of Chern roots generate rationally because the sum of the
roots vanishes on the flag fiber.  Polynomial representatives of a
fiber basis therefore give the required global classes.
The additional degree-two classes have weight two.  Consequently
\[
 \pi^*:\operatorname{Gr}^W_4H^2(\widehat X^{\mathrm B},\mathbb C)
 \xrightarrow{\ \sim\ }
 \operatorname{Gr}^W_4H^2(\widehat Z^{\mathrm B},\mathbb C).
\]
There is a unique class $\xi$ in the space on the left such that
$\iota^*\omega=\pi^*\xi$ on the weight-graded.

By the projection formula, $\mathsf A$ and $\mathsf B$ intertwine
multiplication by $\omega$ and $\xi$ on weight-gradeds.  Their actual
weight shifts are $-2b$ and $2b$.  Equation~\eqref{betti:split} exhibits
the image of $\mathsf B$ as a graded stable summand, with stable
complement $\ker\mathsf A$.  Each Lefschetz isomorphism for $\omega$
is block diagonal for this decomposition and thus restricts to an
isomorphism on the summand.  The half-weight shift of $\mathsf B$ is
$b$, so $\xi$ has center $R+g$ on
$H^*(\widehat X^{\mathrm B},\mathbb C)$.  This correspondence also
makes its cohomology a summand of a Tate twist of the Hodge--Tate
cohomology upstairs, hence Hodge--Tate.

It remains to remove $D_1$.
By Lemma~\ref{norm:low-cohomology} and the given nonabelian-Hodge
diffeomorphism, $H^0(X^{\mathrm B})=\mathbb Q$ and
$H^1(X^{\mathrm B})=0$.  K\"unneth therefore gives
\[
 \xi=\xi_X+\xi_D,
 \qquad
 \xi_X\in\operatorname{Gr}^W_4H^2(X^{\mathrm B},\mathbb C),
 \quad
 \xi_D\in H^2(D_1,\mathbb C).
\]
The tensor-factor Lefschetz lemma, Lemma~\ref{alg:tensor}, shows that
each factor has a common Lefschetz center and that the two centers
add to $R+g$.  The cohomology of $D_1$ is an exterior algebra on $2g$
generators of cohomological degree one and weight two.  Its smallest
and largest half-weight grades are $0$ and $2g$, so its center is
$g$.  Thus $\xi_X$ has center $R$.  Projection and a section also
show that $X^{\mathrm B}$ itself has Hodge--Tate cohomology.

Since $H^0(X^{\mathrm B})\ne0$ and $R>0$, this Lefschetz operator
$\xi_X$ is nonzero.  Lemma~\ref{norm:low-cohomology} gives
$H^2(X^{\mathrm B},\mathbb Q)=\mathbb Q e$ as a vector space, so the
image of $e$ in $\operatorname{Gr}^W_4$ is nonzero and is a nonzero
scalar multiple of $\xi_X$.  Here $e\in W_4H^2$ because
$X^{\mathrm B}$ is smooth.  Rescaling a Lefschetz operator preserves
all its isomorphisms.  They hold over $\mathbb Q$, since the rational
maps in question become isomorphisms over $\mathbb C$.  This proves
Theorem~\ref{betti:lefschetz}.  No step took invariants under the
finite group $\Gamma$.
